% Prose that depends on our editor's results. Numbers come from results.tex (generated); wording is by hand.

\newcommand{\oursDiscussion}{%
\Cref{tab:ours} and \cref{fig:size} show the editor's results. With id addressing, the 1.5B editor is correct on every
one of the 1{,}000 main-tier test edits, and on \editOneFiveHard\,\% of the hard tier, whose networks are larger
than anything in training. On the tasks \astra completed it matches \astra's edits exactly
(\cref{tab:compare}). Because its verdict is computed from its own drawing, its verdicts are correct whenever its
edits are. A 7B editor trained identically reaches the same totals (1{,}000 of 1{,}000 and \editSevenHard\,\%) with
different misses (below). At this task size, model scale buys nothing that the representation does not already
provide.

\paragraph{Addressing matters more than scale.} The same model and data with positional addressing reaches
only \editOneFivePosMain\,\% (circuits \editOneFivePosMainDC\,\%, pipes \editOneFivePosMainPipe\,\%) and
collapses to \editOneFivePosHard\,\% on the hard tier. Accuracy falls steadily with the number of elements in
the drawing (\cref{fig:size}): to name a label as \code{n77} the model must count through the document. In
pipe drawings a pipe's label position follows from its number and the pipe count, while in circuits it
depends on how many sources precede it (each contributes three text elements), and circuits are where
positional addressing fails. Naming elements by their SVG identifier removes the counting and costs nothing at inference.

\paragraph{Opaque-id control.} This advantage depends on how identifiers are named. We replaced every SVG
identifier in each hard-tier source and target with a unique, opaque token, while keeping visible geometry,
labels, instructions and target physics unchanged. All 400 re-derived gold patches still reproduce and solve
their targets. Without retraining, the saved 7B editor falls from 398/400 correct edits to 154/400 (67/200
circuits and 87/200 pipe networks). This is a limit of the editor trained on informative IDs; the verifier
remains identifier-invariant.

\paragraph{Remaining errors and markup.} The 1.5B editor's two hard-tier misses both add a parallel branch to the
largest circuits: the model misplaces the branch's vertical coordinates by 8--30 units, so it no longer taps the
intended wires, on rows lower in the drawing than any seen in training. The 7B editor gets every hard circuit
right. Its two misses add a pipe to the largest water networks, drawing the pipe exactly but placing its label
91--100 units below it, which leaves the new pipe unlabelled. The verifier rejects all four drawings, the first two
because the recovered network differs and the last two because they cannot be read. In the other direction, 22 of
the 1.5B editor's main-tier edits (6 of the 7B's) are correct networks but not byte-identical to the gold drawing.
For the 1.5B editor, 18 leave a junction dot that the gold drawing removes after opening a branch, and 4 place an
added pipe's label 1--12 units from the gold position. A markup metric would
count these as failures. Scoring by the recovered network credits them.}

\newcommand{\costDiscussion}{%
An \astra answer costs \$\astraMainCost{} (main tier) to \$\astraHardCost{} (hard tier) at list price and takes
a median of \astraMainSeconds{}\,s to \astraHardSeconds{}\,s. Our 1.5B editor trains in \trainMinutesOneFivePos{}
minutes on one RTX PRO 6000 GPU and then needs \gpuSecondsPerEditOneFivePos{}\,s of GPU time per edit when batched
sixteen at a time. The id-addressed run shared the GPU with another user's jobs, which stretched its training to
\trainMinutesOneFive{} minutes and inference to \gpuSecondsPerEditOneFive{}\,s per edit. The 7B editor trains in
\trainMinutesSeven{} minutes on one H100 NVL and needs \gpuSecondsPerEditSeven{}\,s per edit. The verifier needs a median of
\verifierMsMain{}\,ms (main) and \verifierMsHard{}\,ms (hard) of laptop CPU time per drawing to recover, solve,
check and compare it. It is cheap enough to run on every edit from any source, including \astra's.}

\newcommand{\conclusionOurs}{A 1.5B editor that names elements by identifier, paired with the verifier, matches
\astra's edits on the tasks both attempted. It is correct on all 1{,}000 main-tier and \editOneFiveHard\,\% of
unseen hard-tier edits, its verdicts are exact by construction, and it runs locally in about a second per edit.}

\newcommand{\abstractOurs}{A 1.5B editor that addresses elements by identifier, paired with the verifier, is correct on
all 1{,}000 held-out edits and \editOneFiveHard\,\% of the larger unseen networks, at about a second of GPU time per edit.}
